Threats a user-level workload confinement deliberately does not address.
Naming them is part of the boundary: an approach honest about what it
does not cover is easier to reason about than one that implies total
coverage. The entries keep their identifiers for citation, and are
grouped where two share one reason.

This set was originally Family 4 of the numbered index, before the
out-of-scope threats were split onto their own \texttt{X} prefix; that
is why the in-scope families run 1, 2, 3, 5 with no Family 4. The
\texttt{X} numbers are what remained.

\textbf{X1 --- a process running as a different real uid}, and
\textbf{X2 --- a process running as root}, are handled by the mechanisms
that already exist for them: ordinary Unix permissions for the first,
root permissions and root-targeted mandatory access control for the
second. A user-level confinement addresses code running as \emph{the
user}, which is the gap those mechanisms leave open; it does not restate
them.

\textbf{X3 --- hardware-level attacks} (cold boot, DMA, side channels,
Spectre-class) and \textbf{X4 --- physical access} are below any
userspace approach's reach. Nothing a reference monitor does in
userspace constrains an attacker with the bus or the hardware, and
claiming otherwise would be the kind of overreach the catalogue exists
to avoid.

\textbf{X5 --- the user knowingly cooperating with the workload.} A user
can grant any capability; an approach surfaces the grant and warns, but
cannot override informed consent, and should not pretend to. This is out
of scope only in its \emph{informed} form --- the user who understands
what a grant permits and makes it anyway. It is distinct from, and must
not be conflated with, the \emph{confused-deputy} case, where the user
is manipulated into carrying an action they did not understand: a
trigger the workload planted in its own tree firing later in the user's
unconfined shell, or a channel the user is tricked into bridging. That
case is squarely in scope and is catalogued at T2.8. The line is
understanding: a deputy who knows is X5, a deputy who is fooled is a
threat the approach is built to close.

\textbf{X6 --- network-level attacks against the host's exposed
services} are a different threat model. A user-level workload
confinement defends the outbound direction --- what the workload reads
and where it connects --- not the host's inbound attack surface, which
is the province of a firewall and the services' own hardening.

\textbf{X7 --- compromise of the underlying kernel primitives the
approach depends on}, and \textbf{X8 --- compromise of the vetted system
tools it depends on}, are beyond the boundary by definition: they are
the trusted substrate the confinement is built on. An approach can keep
that substrate small, pinned, and auditable (and the
confinement-attack-surface family, T5, is about not enlarging it), but a
defect in the kernel's own isolation primitives or in a vetted external
tool is not something the layer above them can contain.

\textbf{X9 --- a confined workload using its legitimate API access to
exfiltrate data} is within-channel exfiltration via an authorised
destination, documented as the T1.8 residual: unblockable at the egress
boundary without destroying the workload's function, and a
data-governance concern rather than an infrastructure one.

\textbf{X10 --- semantic security regressions in workload-produced code
that pattern matching cannot detect} require human review. An approach
can flag known anti-patterns (T2.2), but cannot evaluate code for a
novel logic flaw, an authorization bug, or an injection introduced with
correct syntax and wrong meaning.

\begin{center}\rule{0.5\linewidth}{0.5pt}\end{center}
